\chapter{FUNCTION SPACES}

\textit{July 20th.}

Quizzes 1 and 2 will account for 10\% and 15\% of the final grade respectively. The mid-term will account for 25\% of the final grade, and the final exam will account for 50\%.

\section{An Introduction}
Suppose $\K$ is a field, and $X$ is a topological space. Then the set of all $\K$-valued functions on $X$, denoted by $\K^{X}$, is a vector space over $\K$. If we let $X = \{1,2,\ldots,n\}$, then $\K^{X} \cong \K^{n}$, the $n$-dimensional vector space over $\K$. In this course, we will be primarily concerned with the case when $\K = \R$ or $\C$. Consider the map
\begin{align}
    D:\K^{\{1,2,\ldots,n\}} \to \K^{\{1,2,\ldots,n-1\}}, \quad (Df)(k) = f(k+1)-f(k) \text{ for } 1 \leq k \leq n-1, \; (Df)(n) = f(1)-f(n).
\end{align}
Here, $D$ acts as a discrete analogue of the derivative operator by taking finite differences. We can split this finite difference map as the difference of the left shift operator and the identity operator as $D = S-I$. Since $S$ is the left shift operator, we have $S^{n} = I$ so the eigenvalues of $S$ are a subset of the $n$-th roots of unity. Let $\lambda$ be such an eigenvalue. Then a corresponding eigenvector can be given as $\begin{pmatrix}
    1 \\ \lambda \\ \vdots \\ \lambda^{n-1}
\end{pmatrix}$. If $\lambda = e^{2\pi i /n}$, then we can denote the corresponding eigenvector of $\lambda^{k}$ as $\chi_{k}$. In fact, in the basis $\{\chi_{1},\ldots,\chi_{n-1},\chi_{n}\}$, $D$ is a diagonal matrix. Such a basis is termed a \eax{discrete Fourier basis}, and the coordinate vector of $D$ with respect to this basis is called the \eax{discrete Fourier transform} of $D$. The Fourier transform diagonalizes the finite differential operator. We can view $\chi_{k}$ as a function $\chi_{k}:\{1,2,\ldots,n\} \to \C$ defined by $\chi_{k}(j) = e^{2\pi i k(j-1)/n}$ for $1 \leq j \leq n$. 

Recall the Hermitian inner product on $\C^{n}$ given by $\ip{x,y} = \sum_{i=1}^{n} x_{i} \overline{y_{i}}$. Viewing $\C^{n}$ as the function space $\C^{\{1,2,\ldots,n\}}$, we can endow it with the inner product $\ip{f,g} = \sum_{k=1}^{n} f(k) \overline{g(k)}$, which gives the 2-norm $\norm{f}_{2} = \sqrt{\sum_{i=1}^{n} \abs{f(i)}^{2}}$. The discrete Fourier basis is an orthogonal basis with respect to this inner product. Normalizing via division by $\sqrt{n}$ makes it an orthonormal basis. Thus, if we have $f \in \C^{\{1,2,\ldots,n\}}$ which we wish to write as $f = \sum_{j=1}^{n} a_{j} \chi_{j}$, then we can compute the coefficients $a_{j}$ as
\begin{align}
    a_{j} = \frac{1}{n} \ip{f,\chi_{j}} = \frac{1}{n} \sum_{r=1}^{n} f(r) e^{-2\pi i j(r-1)/n}.
\end{align}
This is the discrete Fourier transform, as before. If we look at the norm, then
\begin{align}
    \norm{f}_{2}^{2} = \norm{\sum_{j=1}^{n} a_{j} \chi_{j}}_{2}^{2} = n \sum_{j=1}^{n} \abs{a_{j}}^{2} \implies \frac{1}{n} \norm{f}_{2}^{2} = \sum_{j=1}^{n} \abs{a_{j}}^{2}.
\end{align}
This is known as \eax{Plancherel's theorem}.

All of the above can be generalized to the continous case; let $L^{2}[0,1]$ be the space of square integrable functions on $[0,1]$. If $(a_{n})_{n \in \Z}$ is the Fourier series of $f$, then $\int_{0}^{1} \abs{f(x)}^{2} dx = \sum_{n \in \Z} \abs{a_{n}}^{2}$, which is the continuous analogue of Plancherel's theorem.

\section{Metric and Normed Spaces}
\textit{July 23rd.}

The following theorem relates the notion of compactness, that every cover has a finite subcover, to many other properties of metric spaces.

\begin{theorem}
Let $(X,d)$ be a metric space. The following are equivalent:
\begin{enumerate}
    \item $X$ is compact.
    \item Every nested sequence of nonempty closed sets $F_1 \supseteq F_2 \supseteq \cdots$ has nonempty intersection, that is,
    \begin{align}
        \bigcap_{n=1}^{\infty} F_n \neq \emptyset.
    \end{align}
    \item Every countable family of closed sets with the finite intersection property has nonempty intersection.
    \item Every sequence in $X$ has a convergent subsequence.
    \item $X$ is complete and totally bounded.
\end{enumerate}
\end{theorem}


\begin{proof}
    For $1 \implies 2$, let $\{F_n\}_{n=1}^{\infty}$ be a nested sequence of nonempty closed sets with $F_1 \supseteq F_2 \supseteq \cdots$. If $\bigcap_{n=1}^{\infty}F_n=\emptyset$, then $\{X\setminus F_n\}_{n=1}^{\infty}$ is an open cover of $X$. By compactness, there exists a finite subcover. Since the $F_n$ are nested,
    \begin{align}
        X=\bigcup_{i=1}^{N}(X\setminus F_i)=X\setminus F_N.
    \end{align}
    Hence $F_N=\emptyset$, a contradiction.

    For $2 \implies 3$, let $\{F_n\}_{n=1}^{\infty}$ be a countable family of closed sets with the finite intersection property. Define
    \begin{align}
        G_n=\bigcap_{i=1}^{n}F_i.
    \end{align}
    Then $G_1\supseteq G_2\supseteq\cdots$ is a nested sequence of nonempty closed sets. By 2, $\bigcap_{n=1}^{\infty}G_n\neq\emptyset$. Since
    \begin{align}
        \bigcap_{n=1}^{\infty}G_n=\bigcap_{n=1}^{\infty}F_n,
    \end{align}
    the result follows.

    For $3 \implies 4$, let $S_0$ be the set of elements in $(x_i)_{i=1}^{\infty}$, a sequence in $X$. Similarly, let $S_1$ be the set of elements in $(x_i)_{i=2}^{\infty}$, and define $S_n$ similarly. Then
    \begin{align}
        \overline{S}_0\supseteq\overline{S}_1\supseteq\cdots\supseteq\overline{S}_n\supseteq\cdots.
    \end{align}
    By 3, $\bigcap_{n=1}^{\infty}\overline{S}_n\neq\emptyset$, so there exists $x'\in\bigcap_{n=1}^{\infty}\overline{S}_n$. Now, $B(x',1)\cap S_1$ cannot be empty, so pick $x_{n_1}$ from this intersection. Again, $B(x',1/2)\cap S_{n_1}$ cannot be empty, so pick $x_{n_2}$ from this intersection. Continue inductively to obtain the $x_{n_k}$'s. One can verify that $x_{n_k}\to x'$.

    For $4 \implies 5$, every Cauchy sequence has a convergent subsequence, and hence converges. Thus, $X$ is complete. Suppose $X$ is not totally bounded. Then there exists $\alpha>0$ such that $X$ has no finite covering by balls of radius $\alpha$. Pick $x_1\in X$. Since $B(x_1,\alpha)\neq X$, pick $x_2\in X\setminus B(x_1,\alpha)$. Again,
    \begin{align}
        B(x_1,\alpha)\cup B(x_2,\alpha)\neq X.
    \end{align}
    Similarly, pick $x_3$ outside the previous balls, and continue inductively. Since $d(x_m,x_n)\ge\alpha$ for $m\neq n$, there is no convergent subsequence, contradicting 4. Hence $X$ is totally bounded.

    For $5 \implies 1$, we prove the contrapositive. Assume $X$ is not compact. If it is not totally bounded, we are done; further assume that it is totally bounded. Then there exists an open cover $\{U_{\alpha}\}_{\alpha\in\Lambda}$ with no finite subcover. Let $\varepsilon=1/2$ and pick a finite $\varepsilon$-net. Pick $x_1$ such that $B(x_1,1/2)$ is not covered by finitely many $U_{\alpha}$. Since $B(x_1,1/2)$ is totally bounded, pick $x_2\in B(x_1,1/2)$ such that $B(x_2,1/4)$ cannot be covered by finitely many $U_{\alpha}$. Inductively, pick $x_n\in B(x_{n-1},1/2^{n-1})$ such that $B(x_n,1/2^n)$ cannot be covered by finitely many $U_{\alpha}$. Then
    \begin{align}
        d(x_n,x_m)\le\sum_{k=m}^{n-1}\frac{1}{2^k},
    \end{align}
    so $(x_n)$ is a Cauchy sequence. Since $X$ is complete, $x_n\to x$ for some $x\in X$. Let $x\in U_{\alpha}$. Then there exists $\varepsilon>0$ such that $B(x,\varepsilon)\subseteq U_{\alpha}$. Moreover, $d(x_n,x)<\varepsilon/2$ for $n$ sufficiently large. Hence
    \begin{align}
        B(x_n,2^{-n})\subseteq B(x_n,\varepsilon/2)\subseteq B(x,\varepsilon)\subseteq U_{\alpha},
    \end{align}
    for all sufficiently large $n$, since $2^{-n}<\varepsilon/2$. This contradicts the choice of $B(x_n,2^{-n})$. Hence $X$ is not complete.
\end{proof}

We now look at some examples of non-trivial function spaces; consider $C[0,1]$, the space of complex-valued continuous functions on $[0,1]$. This is a vector space over $\R$ with pointwise addition and scalar multiplication. We can define a norm on this space as $\norm{f}_{\infty} = \sup_{x \in [0,1]} \abs{f(x)}$. This makes $C[0,1]$ into a normed space.

\begin{definition}
    Let $X$ be a $\K$-vector space. A \eax{norm} on $X$ is a function $\norm{\cdot}:X\to[0,\infty)$ such that for all $x,y\in X$ and $\alpha\in\K$,
    \begin{enumerate}
        \item $\norm{x}=0$ if and only if $x=0$ (positive definiteness).
        \item $\norm{\alpha x}=\abs{\alpha}\norm{x}$ (homogeneity).
        \item $\norm{x+y}\le\norm{x}+\norm{y}$ (triangle inequality).
    \end{enumerate}
\end{definition}

One may verify that $\norm{\cdot}_{\infty}$ is indeed a norm on $C[0,1]$.

\section{The Lebesgue Integral}
\textit{July 27th.}

\begin{definition}
    A function $s:[0,1] \to \C$ is called a \eax{step function} if there exists a partition $P:0=x_{0} < x_{1} < \cdots < x_{n}=1$ of $[0,1]$ such that $s$ is constant on each interval $(x_{i-1},x_{i})$ for $1 \leq i \leq n$. That is, $s(x) = c_{k}$ if $x \in (x_{k-1},x_{k})$ where $c_{k}$ are constants for $1 \leq k \leq n$.
\end{definition}

We impose that the integral of such a step function is given by
\begin{align}
    \int_{0}^{1} s(x) \, dx = \sum_{k=1}^{n} c_{k}(x_{k}-x_{k-1}).
\end{align}
The values of the step function on the partition points $x_{i}$ do not affect the value of the integral.

\begin{proposition}
    The set of step functions on $[0,1]$ is a complex vector space under pointwise addition and scalar multiplication.
\end{proposition}
The proof is straightforward and left as an exercise. 

\begin{proposition}
    The set of step functions on $[0,1]$ such that $\int_{0}^{1} \abs{s(x)} \, dx = 0$ is a subspace of the vector space of step functions on $[0,1]$.
\end{proposition}
Again, this proof is rudimentary and left as an exercise. The reader is also urged to show that $\int_{0}^{1} s(x) \, dx$ is well-defined and independent of the choice of partition $P$.

\begin{definition}
    A set $S \subseteq [0,1]$ is said to have \eax{content zero} or measure zero if for every $\eps > 0$, there exists countably many open intervals $(a_{n},b_{n})_{n=1}^{\infty}$ such that $S \subseteq \bigcup_{n=1}^{\infty} (a_{n},b_{n})$ and $\sum_{n=1}^{\infty} (b_{n}-a_{n}) < \eps$.
\end{definition}

As a trivial example, any finite subset of $[0,1]$ has content zero. Similarly, any countable subset of $[0,1]$ has content zero. The Cantor set is an example of an uncountable set with content zero. One may show that the countable union of sets of content zero has content zero.

\begin{definition}
    A sequence of functions $f_{n}:[0,1] \to \R$ is said to be increasing if $f_{n}(x) \leq f_{n+1}(x)$ for all $x \in [0,1]$. Similarly, a sequence of decreasing functions is defined.
\end{definition}

Note that a property is said to be hold \eax{almost everywhere} if it holds for all $x \in [0,1]$ except for a negligible set (a set of content zero).

\begin{definition}
    We say $f_{n} \to f$ pointwise almost everywhere if there exists a set $S \subseteq [0,1]$ of content zero such that $f_{n}(x) \to f(x)$ for all $x \in [0,1] \setminus S$. That is, $f_{n}$ converges to $f$ pointwise except on a negligible set.
\end{definition}

\begin{theorem}[The \eax{proto monotone convergence theorem}]
    Let $\{s_{n}\}_{n=1}^{\infty}$ be a decreasing seqeucence of non-negative step functions on $[0,1]$ such that $s_{n} \to 0$ pointwise almost everywhere. Then $\int_{0}^{1} s_{n}(x) \, dx \to 0$ as $n \to \infty$.
\end{theorem}

\begin{proof}
    Let $E \defeq \{x \in [0,1] \mid s_{n}(x) \not\to 0\}$. Then $E$ has content zero. Define $D$ to be the set of endpoints of intervals in the partitions corresponding to the $s_{n}$'s. Then $D$ is countable, and hence has content zero. Let $F = D \cup E$ is negligible. Let $x \in [0,1]\setminus F$. Then there exists $N(x) \in \N$ such that $s_{n}(x) = 0$ for all $n \geq N(x)$. Since $x \notin D$, there exists an open interval $B(x)$ such that $s_{N}(t) < \eps$ for all $t \in B(x)$. As $\{s_{n}\}$ is decreasing, we have $s_{n}(t) < \eps$ for all $t \in B(x)$ and $n \geq N(x)$. For every $x \in [0,1]\setminus F$, we have a $B(x)$. Find a sequence of open intervals such that $F \subseteq \bigcup_{n=1}^{\infty} F_{n}$ and $\sum_{n=1}^{\infty} \abs{F_{n}} < \eps$. Then
    \begin{align}
        [0,1] \subseteq \left( \bigcup_{x \in [0,1]\setminus F} B(x) \right) \cup \left( \bigcup_{n=1}^{\infty} F_{n} \right),
    \end{align}
    an open cover. By compactness, there exists a finite subcover, so
    \begin{align}
        [0,1] \subseteq \bigcup_{i=1}^{p} B(x_{i}) \cup \bigcup_{r=1}^{q} F_{r}.
    \end{align}
    Call the left union $P$ and the right union $Q$. Then
    \begin{align}
        \int_{Q} s_{n}(x) \, dx \defeq \int_{0}^{1} \chi_{Q}(x) s_{n}(x) \, dx \leq M \int \sum_{i=1}^{q} \abs{F_{i}} \, dx < M \eps.
    \end{align}
    Similarly,
    \begin{align}
        \int_{P} s_{n}(x) \, dx \defeq \int_{0}^{1} \chi_{P}(x) s_{n}(x) \, dx \leq \eps \quad \text{ for all } n \geq \max\{N(x_{i}) \mid 1 \leq i \leq p\}.
    \end{align}
    Combining gives
    \begin{align}
        \int_{0}^{1} s_{n}(x) \, dx = \int_{P} s_{n}(x) \, dx + \int_{Q} s_{n}(x) \, dx < (M+1)\eps.
    \end{align}
    Thus, $\int_{0}^{1} s_{n}(x) \, dx \to 0$ as $n \to \infty$.
\end{proof}

From here, we can immediately note that $\rho:\cS[0,1] \to \C$, map from the set of step functions, defined by $s \mapsto \int_{0}^{1} s(x) \, dx$ is a linear functional. We can also define a semi-norm on the set of step functions as $\norm{s}_{1} = \int_{0}^{1} \abs{s(x)} \, dx$. This is a semi-norm since $\norm{s}_{1} = 0$ does not imply that $s = 0$; it only implies that $s = 0$ almost everywhere. We can define an equivalence relation on the set of step functions by $s \sim t$ if and only if $s = t$ almost everywhere. Then we can define a norm on the quotient space $\cS[0,1]/\sim$ as $\norm{s}_{1} = \int_{0}^{1} \abs{s(x)} \, dx$. This is indeed a norm since $\norm{s}_{1} = 0$ implies that $s = 0$ in the quotient space.

The normed space $(\cS[0,1]/\sim, \norm{\cdot}_{1})$ is not complete; there is a much larger class of functions that are (Lebesgue) integrable. Henceforth, we define $L^{1}[0,1]$ to be the completion of $(\cS[0,1]/\sim, \norm{\cdot}_{1})$. The elements of $L^{1}[0,1]$ are called \eax{Lebesgue integrable function}s. 